% ERL Letter: Methods, concise (revised after the internal review). Full
% detail is in the Supplementary Information (supplementary_erl.tex, section \SImethods; the
% revision analyses in \SIrevision). Every number is a \metric{} macro. Equations (1)-(3) were
% checked against the functions named in the comment beside each.

\subsection{Observed shorelines}
The study domain is the North Caspian north of the line from Chechen Island to Cape Tyub-Karagan (\metric{study_area.domain_km2}~\kmsq{}, land included), divided before any evaluation into seven sectors: the Kizlyar Bay--Dagestan--Kalmyk coast (west), the Volga and Ural deltas, the Volga--Ural interfluve, the northeastern Kazakh coast, Komsomolets Bay and the open shelf.

We used the Landsat~5, 7, 8 and 9 Collection~2 Level-2 scenes over the domain (\metric{erl.text.inventory.n_landsat} scenes)\citep{usgs_l45tm_c2l2,usgs_l7etm_c2l2,usgs_l89oli_c2l2}, processed in Google Earth Engine\citep{gorelick2017}. Clouds, shadows, saturated pixels and the ice months (\metric{params.ice_first_month} to \metric{params.ice_last_month}) were masked. Each pixel was classified as water or dry with the modified normalised difference water index (MNDWI)\citep{xu2006} and a threshold set once per acquisition from its histogram; pixels with a normalised difference vegetation index above \metric{params.ndvi_veg} are emergent reeds, which hide the water's edge, and were treated as unobserved.

Each scene was given the basin level at its acquisition time from the altimetry series of the Database for Hydrological Time Series of Inland Waters (DAHITI)\citep{schwatke2015}, \metric{levels.dahiti.n_records} records from \metric{levels.dahiti.first} to \metric{levels.dahiti.last}, with long gaps filled from Hydroweb\citep{hydrowebwl,cretaux2011} (median 1$\sigma$ \metric{levels.scene_sigma_p50_m}~m); gauge posts\citep{caspcom2026} tie it to the Baltic Height System.

Scenes were grouped into level bins \metric{params.bin_width_m}~m wide, separately for the rising limb (E1, to \metric{levels.epoch_E1_end}), the early fall (E2, to \metric{params.e2_end}) and the late fall (E3). Per \metric{params.stage2_res_m}~m cell, a level's observed sea consists of the cells that were water in at least \metric{params.contour_level} of at least \metric{params.min_valid_obs} valid acquisitions and connect to the open sea. Cells that are sea at the record low were kept as sea at every higher level, which removes deep, clear water misread as land. Alternate acquisitions of each bin form two interleaved halves, whose shorelines give the observation noise floor. The floor measures precision, not errors common to both halves, and at levels crossed on both limbs it also holds distances between epochs, as registered.

Four sectors are validated (west, interfluve, northeastern coast, open shelf): there the waterline elevation model built from the same observations (below) agrees with heights of the Ice, Cloud and land Elevation Satellite-2 (ICESat-2). It is biased in the Volga delta and Komsomolets Bay, and the Ural delta has no test block.

\subsection{Hindcast audit}
We audited five global grids: GEBCO\_2026\citep{gebco2026}, GEBCO\_2022\citep{gebco2022}, ETOPO 2022\citep{etopo2022,macferrin2025}, ETOPO1\citep{amante2009} and SRTM15+ V2.7\citep{tozer2019}. They share sources, so they are partly dependent\citep{macferrin2025,gebco2022}. With $z_j$ a grid's height of cell $j$ on the DAHITI reference and $S$ the open-sea seed cells, the flood level of cell $i$ is
% src/caspian/hindcast_tiled.py: flood_level_map (minimax path height, 8-connected, finite cells);
% src/caspian/products.py: flood_level (threshold L + c); scripts/a7_hindcast.py: fit_offset (c)
\begin{equation}
  F_i = \min_{\pi \in \Pi(S, i)} \; \max_{j \in \pi} z_j ,
  \label{eq:floodlevel}
\end{equation}
where $\Pi(S, i)$ is the set of 8-connected paths of cells with finite heights from $S$ to $i$; cell $i$ is predicted sea at level $L$ when $F_i < L + c$. The grids' vertical reference in the Caspian is uncertain\citep{gebco2026,tozer2019}, so $c$ is the single constant per grid that minimises its summed absolute area error over all observed levels on training blocks (the shape variant, which decides); it is best-fitting for area, not for shoreline position. $c = 0$ on the documented datum chain is the total variant. Blocks of \metric{params.block_size_km}~km were assigned at random to training, calibration and test in proportions \metric{params.split_train}, \metric{params.split_calib} and \metric{params.split_test}, stratified by sector (seed \metric{params.seed}).

The displacement $D_{s,L}$ of a grid in sector $s$ at level $L$ is the median distance between its predicted shoreline cells and the observed ones, pooled in both directions and over epochs. With $Q_{s,L}$ the \metric{params.err_pct_hi}th percentile of the same distances between the two observation halves, the grid mispredicts sector $s$ if
% src/caspian/h1.py: mispredicts (exceeds = median > p95_noise; mispredicts = n_exceeding > n_evaluable / 2)
\begin{equation}
  \sum_{L \in \mathcal{L}_s} \mathbf{1}\!\left[ D_{s,L} > Q_{s,L} \right] > \tfrac{1}{2}\,\lvert \mathcal{L}_s \rvert ,
  \label{eq:mispredicts}
\end{equation}
where $\mathcal{L}_s$ is the set of evaluable levels of the sector (with an observed shoreline, a grid distance and a noise distance). We report the per-level ratio $D/Q$ and the probability of each verdict under a moving-block bootstrap over levels (blocks of \metric{rev.w.s.boot.block_m}~m, \metric{rev.w.s.boot.n} resamples). Larger errors in the deltaic and vegetated sectors (west, Volga, Ural) were tested with a one-sided permutation test of block medians (\metric{params.permutations_n} permutations, $\alpha = \metric{params.alpha}$), and after review at sector level with Holm's adjustment.

\subsection{Great Salt Lake}
At Great Salt Lake the same pipeline audited the north and south arms, with gauge levels of the United States Geological Survey (USGS) on the National Geodetic Vertical Datum of 1929\citep{usgs_nwis}, converted by us to the North American Vertical Datum of 1988 with VERTCON~3.0\citep{noaa_ncat}. The registered control, a USGS topobathymetric model\citep{root2023_gsl_tbdem}, could not be retrieved by script, so its registered substitute served: the 3D Elevation Program (3DEP) model\citep{usgs_3dep10m}, built over the lake bed mostly from lidar of autumn 2016\citep{quantumspatial2017_gsl,usgs_3dep13_tiles}, hydro-flattened\citep{archuleta2017} and evaluated only above its flattened surfaces.

\subsection{Near-term emergence}
The 2026 reference level $L_{2026}$ is \metric{h3.level_2026_m}~m on the DAHITI reference, the January--September mean of the \metric{h3.level_2026.n_records} altimetry records of 2026. For a decline $\Delta$ of \metric{h3.scenarios_text}~m, a grid's emerged area referenced to its own 2026 shoreline is
% src/caspian/exposure.py: emerged_curve; scripts/a8_emergence.py (sea cells F < L in the domain)
\begin{equation}
  E(\Delta) = A(L_{2026}) - A(L_{2026} - \Delta) ,
  \label{eq:emerged}
\end{equation}
where $A(L)$ is the area of the domain cells with $F_i < L$. The grids as distributed ($c = 0$, the registered primary) only illustrate datum risk. We report $E(\Delta)$ with each grid shifted by its constant $c$, and from the observed 2026 shoreline: the observed sea at the lowest audited level, completed with the 2025 and 2026 near-minimum extent, less the grid's sea at $L_{2026} - \Delta$. This estimates the seabed under the 2026 sea that a uniform fall $\Delta$ would expose, with the seabed's shape from the grid; it assumes a static seabed and a flat water surface and inherits the grid's shape error. Hindcast errors of emerged area between observed levels, normalised by the decline, give error envelopes that describe past errors, not coverage, up to the largest hindcast decline. Simulated correlated vertical errors at each grid's ICESat-2 dispersion (correlation ranges of \metric{rev.w.n.sim.range_secondary_km} and \metric{rev.w.n.sim.range_primary_km}~km) give a reference for agreement between grids. Exposure (OpenStreetMap\citep{osm}, ESA WorldCover\citep{worldcover2021}, protected areas\citep{wdpa}) is a diagnostic without intervals. Pathways\citep{samant2023,hoseini2025} are applied from the 2026 level, and emerged area at fixed years is split between pathways and grids.

\subsection{Waterline elevation model and independent heights}
We interpolated the E2 and E3 waterlines into a \metric{params.dem_res_m}~m elevation model of the emerged band and tested it, as registered, against ICESat-2 heights\citep{atl06_v7,atl08_v7,markus2017} over seabed already dry when measured, on held-out \metric{h2.block_size_km}~km blocks, with the Copernicus 30~m elevation model\citep{copdem,copdem_handbook_i5} and the Global Ecosystem Dynamics Investigation (GEDI)\citep{gedi_l2a_v2} as further checks (Supplementary section~\SIdem).

\subsection{Pre-registration and reproducibility}
The hypotheses, metrics, blocks and decision rules were registered before any hindcast, topography-accuracy or projection metric was computed (registration and amendments: Supplementary file; deviations: section~\SIdeviations). Confirmatory claims are the registered verdicts (H1 by equation~(\ref{eq:mispredicts}) in the seven sectors, H2 and H3) and their registered sensitivities. The analyses added after internal review, including the validated-sector headline, were registered as Amendment 18 before they were computed and are exploratory; Supplementary section~\SIrevision{} classifies every quantitative claim.
